%% ============================================================
%% FIGURE INSERT INSTRUCTIONS FOR YOUR OVERLEAF PAPER
%% Upload these 5 images to your Overleaf images/ folder:
%%   reno_cwnd_panels.png
%%   reno_throughput_cdf.png
%%   reno_metrics_comparison.png
%%   cubic_cwnd_panels.png
%%   cubic_metrics_comparison.png
%% ============================================================


%% ============================================================
%% INSERT #1: AFTER Table V (tab:reno) in Section IV-B
%% (After the paragraph starting "The proportional mechanism 
%%  proactively injects 18 ECN signals...")
%% ============================================================

\begin{figure}[!t]
\centering
\includegraphics[width=\columnwidth]{reno_cwnd_panels.png}
\caption{TCP Reno CWND traces under three conditions (120\,s, 10\,Mbit/s, 30\,ms RTT). Top: Baseline---frequent deep drops to 1--2 segments. Middle: Binary ECN---similar oscillation pattern with no active signals fired. Bottom: Proportional ECN---higher sustained CWND ($\sim$6.0 segments) with 18 proactive ECN signals maintaining utilization closer to bottleneck capacity.}
\label{fig:reno_cwnd}
\end{figure}

\begin{figure}[!t]
\centering
\includegraphics[width=\columnwidth]{reno_throughput_cdf.png}
\caption{Empirical CDF of instantaneous throughput for TCP Reno. The proportional ECN curve (green) is consistently right-shifted, indicating higher throughput across all percentiles. At the median, our approach achieves $\sim$1.5\,Mbps vs.\ $\sim$1.0\,Mbps for baseline---a 15.5\% improvement in average throughput.}
\label{fig:reno_cdf}
\end{figure}

\begin{figure}[!t]
\centering
\includegraphics[width=\columnwidth]{reno_metrics_comparison.png}
\caption{TCP Reno performance metrics comparison. Left: throughput progression (1.39$\rightarrow$1.49$\rightarrow$1.60\,Mbps). Middle: retransmission counts from \texttt{ss} snapshots; note that iperf3 recorded 1,077 cumulative retransmits for baseline vs.\ 0 for both ECN conditions. Right: CWND stability ($\sigma$).}
\label{fig:reno_metrics}
\end{figure}


%% ============================================================
%% INSERT #2: AFTER Table VI (tab:cubic) in Section IV-C  
%% (After the paragraph starting "Our proportional approach 
%%  achieves the best CWND stability...")
%% ============================================================

\begin{figure}[!t]
\centering
\includegraphics[width=\columnwidth]{cubic_cwnd_panels.png}
\caption{TCP Cubic CWND traces under three conditions. Top: Baseline---large oscillations reaching 12--13 segments with deep drops, causing 755 retransmissions. Middle: Binary ECN---narrower oscillation band with zero retransmissions. Bottom: Proportional ECN---smoothest trace ($\sigma=1.80$, an 11.3\% stability improvement over baseline) with 2 targeted ECN signals at peak-confidence prediction points.}
\label{fig:cubic_cwnd}
\end{figure}

\begin{figure}[!t]
\centering
\includegraphics[width=\columnwidth]{cubic_metrics_comparison.png}
\caption{TCP Cubic performance metrics comparison. The retransmission panel provides the strongest visual result: 755 baseline retransmissions reduced to exactly 0 by both ECN approaches. The modest throughput reduction ($-9.4\%$) represents the well-documented ECN trade-off: preventing queue overflow keeps CWND below saturation, accepting marginally lower throughput in exchange for zero packet loss and 11.3\% CWND stability improvement.}
\label{fig:cubic_metrics}
\end{figure}
